\chapter{Introducción}
\label{ch:introduccion}

\section{Motivación y contexto}

La gestión cuantitativa de inversiones ha cambiado mucho en las últimas
décadas. El modelo de selección de carteras de Harry Markowitz (1952) puso las
bases matemáticas de la diversificación y del equilibrio entre rentabilidad y
riesgo. En la práctica, sin embargo, choca con un problema: hay que estimar sus
parámetros de entrada, y sobre todo las rentabilidades esperadas de los activos.

\section{Planteamiento del problema}

El modelo de Markowitz necesita como entradas el vector de rentabilidades
esperadas $\boldsymbol{\mu}$ y la matriz de covarianza $\boldsymbol{\Sigma}$. La
matriz $\boldsymbol{\Sigma}$ se estima con cierta estabilidad a partir de datos
históricos, y todavía mejora con técnicas de shrinkage. Estimar
$\boldsymbol{\mu}$, en cambio, es mucho más difícil: pequeños errores en la
predicción de rentabilidades dan carteras muy distintas y, a menudo, peores que
un simple reparto equiponderado \cite{DeMiguel2009}.

Este trabajo estudia si el Machine Learning puede mejorar la predicción de
$\boldsymbol{\mu}$ y, con ello, el rendimiento \emph{out-of-sample} (fuera de la
muestra de estimación) de las carteras optimizadas.

\section{Objetivos}

\subsection{Objetivo general}

Diseñar y evaluar un modelo de optimización de carteras que integre la
teoría clásica de Markowitz con técnicas de Machine Learning para la
predicción de rentabilidades esperadas.

\subsection{Objetivos específicos}

\begin{enumerate}
    \item Formular el modelo media-varianza como un problema cuadrático
          convexo (QP) y resolverlo de forma exacta mediante \texttt{cvxpy}.
    \item Implementar y comparar tres estimadores de la matriz de covarianza:
          muestral, shrinkage de Ledoit-Wolf y basado en PCA/factores.
    \item Desarrollar modelos de Machine Learning (Ridge, Lasso, Random Forest)
          para predecir las rentabilidades esperadas de cada activo.
    \item Evaluar el rendimiento out-of-sample de las carteras mediante
          backtesting walk-forward, comparándolas contra baselines relevantes.
\end{enumerate}

\section{Estructura del documento}

El resto del documento se organiza como sigue. El Capítulo~\ref{ch:estado-arte}
revisa el estado del arte en optimización de carteras, desde la teoría moderna
hasta el aprendizaje automático. El Capítulo~\ref{ch:marco-teorico}
presenta el marco teórico financiero. El Capítulo~\ref{ch:formulacion} desarrolla
la formulación matemática del problema como QP convexo, incluyendo las
condiciones KKT y la solución analítica de la frontera eficiente. El
Capítulo~\ref{ch:limitaciones} aborda las limitaciones del modelo clásico
y presenta los estimadores robustos de covarianza. El Capítulo~\ref{ch:ml}
describe la aplicación de Machine Learning a la predicción de rentabilidades.
El Capítulo~\ref{ch:metodologia} detalla la metodología experimental. El
Capítulo~\ref{ch:resultados} presenta y discute los resultados. Finalmente,
el Capítulo~\ref{ch:conclusiones} recoge las conclusiones y posibles líneas
futuras de investigación.
